\section{Resultados}
\label{sec:resultados}

% ==== Opción 1: pega aquí los resultados generados por el Cuaderno 04 ====
% El Cuaderno 04 exporta un archivo `resultados.tex` (matriz, vector y tabla de
% métodos). Cópialo a esta carpeta y descomenta la siguiente línea:
% \input{resultados}

\subsection{Solución del sistema}
La solución obtenida por los distintos métodos fue:

\begin{table}[h]
	\centering
	\begin{tabular}{lc}
		\toprule
		Método & Solución $\mathbf{x}$ \\
		\midrule
		Eliminación Gaussiana & (escribe aquí) \\
		Descomposición LU     & (escribe aquí) \\
		Gauss--Seidel         & (escribe aquí) \\
		Thomas / Cholesky     & (escribe aquí) \\
		\texttt{np.linalg.solve} (referencia) & (escribe aquí) \\
		\bottomrule
	\end{tabular}
	\caption{Comparación de la solución por método.}
	\label{tab:metodos}
\end{table}

\subsection{Análisis de Monte Carlo}
Se perturbó el vector $\mathbf{b}$ con ruido relativo y se resolvió el sistema $N$ veces
para evaluar la sensibilidad de la solución. La figura~\ref{fig:montecarlo} muestra la
distribución resultante.

\begin{figure}[h]
	\centering
	% \includegraphics[width=0.8\textwidth]{img/mc_boxplot.png}
	\caption{Box plot de la solución ante incertidumbre en $\mathbf{b}$.}
	\label{fig:montecarlo}
\end{figure}

% Comenta: ¿el sistema está bien o mal condicionado? ¿Qué implica el valor de Cond(A)?
